\chapter{KIỂM THỬ, ĐÁNH GIÁ VÀ KẾT LUẬN}

\section{Kiểm thử chức năng}

Các tình huống kiểm thử chính được thực hiện trên giao diện và logic trò chơi như sau:

\begin{table}[H]
\centering
\caption{Bảng kết quả kiểm thử chức năng}
\label{tab:kiem-thu-chuc-nang}
\begin{tabular}{|p{4cm}|p{5cm}|p{4cm}|}
\hline
\textbf{Test case} & \textbf{Hành động} & \textbf{Kết quả} \
\hline
Khởi động trò chơi & Mở trang chủ và bấm nút Start & Trò chơi vào đúng màn chơi ban đầu. \
\hline
Lưu tiến trình & Bấm nút lưu tiến trình trong lúc chơi & Trạng thái được lưu lại vào localStorage. \
\hline
Giải câu đố đúng & Nhập đáp án đúng cho câu đố hiện tại & Điểm tăng và gem bị tiêu thụ. \
\hline
Giải câu đố sai & Nhập đáp án sai nhiều lần & Hệ thống trả về phản hồi và cho phép thử lại. \
\hline
Bảng xếp hạng & Hoàn thành trò chơi và nhập tên & Điểm được thêm vào bảng xếp hạng. \
\hline
\end{tabular}
\Nguon{Kết quả kiểm thử cục bộ trên giao diện trò chơi}
\end{table}

\section{Kiểm thử bảo mật}

Một số tình huống bảo mật đã được xem xét như sau:

\begin{itemize}
    \item Test dữ liệu bị sửa: thay đổi giá trị điểm trong localStorage để xem liệu hệ thống có bị vượt qua giới hạn không.
    \item Test sai khóa: nhập khóa sai cho các câu đố Vigenère hoặc RSA để quan sát phản hồi.
    \item Test sai chữ ký hoặc sai HMAC: trong mô hình minh họa, dữ liệu bị thay đổi sẽ không được chấp nhận khi kiểm tra không đúng.
    \item Test replay: lặp lại cùng một hành động nhiều lần và nhận thấy điểm số hoặc tiến trình cần được kiểm soát chặt chẽ hơn.
    \item Test quyền: các file dữ liệu và khóa đều được quản lý tập trung trong thư mục data, nên việc truy cập không được kiểm soát ở mức sản phẩm hoàn chỉnh.
\end{itemize}

Nhận xét: trò chơi có giá trị minh họa tốt cho các khái niệm bảo mật, nhưng vẫn cần thêm cơ chế xác thực phía server và kiểm soát dữ liệu nghiêm ngặt hơn khi triển khai thực tế.

\section{Benchmark hiệu năng}

Để đánh giá hiệu năng, nhóm quan sát thời gian xử lý ở mức mô phỏng trên máy cục bộ:

\begin{table}[H]
\centering
\caption{Benchmark hiệu năng cục bộ}
\label{tab:benchmark}
\begin{tabular}{|p{4cm}|p{4cm}|p{4cm}|}
\hline
\textbf{Tác vụ} & \textbf{Kích thước / số lượng} & \textbf{Kết quả quan sát} \
\hline
Tạo câu đố Caesar & 10 câu & Thời gian xử lý rất nhanh. \
\hline
Tạo câu đố Vigenère & 10 câu & Thời gian xử lý nhanh, ổn định. \
\hline
Giải mã RSA minh họa & 1 câu & Thời gian chấp nhận được trong môi trường local. \
\hline
Lưu tiến trình & 1 lần & Hoạt động tức thời trên trình duyệt. \
\hline
\end{tabular}
\Nguon{Thử nghiệm cục bộ trên máy phát triển}
\end{table}

\section{Bảng kế thừa và nâng cấp}

Phần nội dung này kế thừa ý tưởng từ một số mẫu game giáo dục trước đó về việc dùng trò chơi để minh họa các khái niệm mật mã. Nhóm tự nâng cấp bằng cách bổ sung các tác nhân tấn công giả định, các câu đố liên quan đến lỗ hổng bảo mật, bảng xếp hạng, lưu tiến trình và thiết kế giao diện gần với môi trường học tập thực tế.

\section{Kết luận và hướng phát triển}

Dự án đã hoàn thành được mục tiêu ban đầu là xây dựng một trò chơi giáo dục giúp người học hiểu được các nguyên lý bảo mật thông qua trải nghiệm trực quan. Tuy nhiên, hệ thống vẫn còn hạn chế ở chỗ chưa có lớp bảo mật phía server, chưa lưu trữ dữ liệu an toàn bằng cơ sở dữ liệu và chưa tích hợp xác thực người dùng hoặc kiểm soát phiên. Trong tương lai, nhóm có thể nâng cấp bằng cách thêm cơ chế xác thực, lưu trữ dữ liệu trên server, triển khai HTTPS, và mở rộng các bài học về bảo mật hiện đại như TLS, JWT, OAuth và bảo mật ứng dụng web.